% Prospective apparatus and protocol; scientific and publication acceptance remain open.
\section{Build a brain--body system we can test}

What would justify building the next prototype? The paired observations in
\emph{From a reach to a testable model} and the tasks in \emph{Catch it again
when the help is gone} give that decision a test. The virtual hand would predict
from one declared input side; independent movement observations would check the
answer. A missed contact or a late movement identifies a problem to resolve
before expanding the apparatus.

The proposed design brings non-invasive EEG and surface EMG into one conditional
learner with the calibrated motion-capture relay-composite and an open
musculoskeletal-simulation physics-engine shown in a lifelike view. Its open
multimodal model would need eligible neural and muscular structure and activity
data. The EEG/EMG pair supplies intent-related brain and muscle observations;
the movement/physics pair tests what the body did and how the mechanics predict
it. Form constrains what could happen; measured function checks the prediction.
Neither pair can certify itself through agreement inside the learner.

These five roles give us brain recordings, muscle recordings, measured
movement, a computer body and a learner that relates them. The capture
composite would join checked movement views. The physics engine would
calculate movement under stated muscle, joint and contact rules. Multimodal
means that the learner can relate different kinds of observation; it does
not mean they were all recorded together or make a missing channel measured.

\subsection{One movement. Two directions.}

The two directional tests need different answers kept outside their inputs.
For the forward test, EEG and declared task context would predict supported
muscle activity and movement on withheld recordings. For the reverse test,
muscle activity and movement would predict recorded EEG features within the
task. Each prediction needs its own independent measured target; agreement
between models cannot substitute for either comparison.

Movement establishes what the body did, narrowing the plausible explanations
of the brain signals without identifying a unique intent or private experience.
Stimulus, task, behaviour and separately reported context therefore remain with
the observations. Uncertainty belongs beside the estimate; a generated neural
pattern cannot be displayed as the observed mind.

Surface EMG still records muscle electrical activity. An incoming sensory or
feedback-pathway claim needs its own evidence, which the virtual body cannot
supply. Recordings, predictions and simulation must remain identifiable as the
comparison follows one direction and then the other.

\subsection{A bridge must account for the spine}

The spinal cord helps shape the response. Muscle output depends on sensory
input, movement context, local circuitry and descending input; some spinal
pathways can retain learned changes. Reproducing a specified spinal function
would therefore require more than forwarding its signal: a brain--body bridge
must account for the processing it aims to reproduce.

A changed response can reflect immediate reflex modulation, lasting spinal
plasticity or peripheral adaptation. Those explanations need separate tests.
A change during the task does not establish retained memory, and a retained
change does not by itself locate that memory. Human reflex conditioning and
constrained mouse adaptation supply evidence for different claims.~\cite{thompson2009retained,thompson2013retained,lavaud2024}
The human studies do not establish brain-independent learning; the mouse
preparation does not demonstrate unrestricted walking or reaching. Our human
comparisons use abstracts and selected results, discussion or methods passages.
Relevant mouse-study passages were also examined in the author manuscript;
equivalence to the final published edition and the complete methods and
statistics have not been established here.

Reflex modulation means that the current conditions change a response;
plasticity means a change in the nervous system that can last. Peripheral
adaptation concerns changes outside the brain and spinal cord. A later
change in movement could involve more than one of these processes, which
is why persistence alone cannot tell us where or how the change occurred.

The periphery adds a different question: did the sensory signal itself change?
Passive proprioceptive training supplies a task-specific peripheral comparison;
its psychophysical and spindle-afferent experiments were separate.~\cite{ackerley2022retained}
Our passive-training comparison uses abstracts and selected results, discussion
or methods passages. None of these studies demonstrates universal skill
transfer or complex memory stored in a peripheral axon.

The psychophysical comparison tests what a person can distinguish or report.
The spindle-afferent comparison records sensory signals associated with
muscle spindles, receptors involved in sensing stretch. These were separate
experiments: placing their results together does not turn them into a single
simultaneous observation of perception and nerve activity.

The next comparison must choose a function we can measure. Does the controller
respond appropriately to the supported context? Does an adaptation persist
through the declared test? Could peripheral or mechanical changes explain the
outcome? Keep recordings, estimates and simulated state in distinct visible
layers. Surface EMG cannot supply a missing sensory-nerve or spinal channel;
an explanatory cutaway cannot register recorded neural sites.

Externalization remains the research aim: can a specified spinal function be
reproduced outside its usual pathway? Replacing a computation, stimulating a
pathway, controlling a simulated limb and restoring biological function would
answer different questions. Each needs its own task- and site-specific evidence;
no programme spinal intervention or qualified spinal-writing interface is reported. Anatomical access does not establish
neural function, increased excitability does not establish useful plasticity,
and no fixed percentage of a complex skill can be allocated to brain, spine
or periphery from these proposals.

\subsection{Reading a signal and writing an output}

A prediction leaves the delivery question open: what, if anything, reached
the trainee? \emph{Write sensation. Enact action} distinguishes the computer's
input/output boundary from afferent/efferent direction relative to the central
nervous system; a robot command has an action role without becoming an
anatomical efferent nerve. That distinction determines the further test.

Delivering sensation would require a separately qualified sensory writer or
transducer. Indirect sensory feedback and direct neural writing need distinct
device, measurement, safety and clinical evidence. Producing an action needs
the appropriate virtual, robotic or biological output interface. The recording
prototype does neither by predicting EEG or EMG. Physical sensory writers,
transducers and body actuators remain excluded from the acquisition and
virtual-work plan, unqualified and unpriced. No stimulation, participant acquisition, biological actuation or training through these physical interfaces has been performed or qualified here.

\subsection{The first prototype has to earn the next one}

The lean tier would first ask whether the selected signals and calibrated
movement reference can be synchronized for a bounded task. The serious-prototype
tier would improve acquisition, independent references and source-specific
bidirectional evaluation. Evidence from those comparisons would determine
whether the research-lab tier should expand cohorts, tasks, instrumentation and
compute. A larger apparatus needs an answered question from the smaller one.

Equipment, calibration, accessories, storage, integration, independent references, safety and consent work, and contingency would all affect the cost. No complete, dated cost estimate has been established for these tiers. The lean, serious-prototype and research-lab scopes therefore specify different experiments, not demonstrated price points or a complete neural reconstruction. The unfulfilled currency, vendor and quotation requirements remain in the budget-preparation provenance record.

\subsection{Let the model earn the claim}

The first model comparison depends on compatible observations, not the
length of the source catalogue. Rights, version, coverage, clock, units and
missing modalities determine what an open dataset can contribute. Fitting
would require an independently admitted protocol, with participants, sessions
and tasks excluded from model selection reserved for testing. Simple baselines
would show whether the added model helps; artifact leakage, absent inputs,
timing errors, physically invalid solutions and uncertainty could explain where
it fails. Earlier failures and reserved recordings remain part of that test's
constraints.

Constrained virtual hand or foot control would be the starting domain. Robotics,
remote operation, unfamiliar effectors and extra limbs remain later research
directions, each requiring separate physical, latency and safety evidence.
A successful movement prediction supports its movement claim. Biological and
perceptual claims need the corresponding broader evidence.

Rights-cleared images and lifelike interactive models would let a reader inspect
the apparatus and compare the observation, prediction and error under their
declared conditions. A convincing image cannot fill a missing measurement or
demonstrate biological function. The first prototype would have to make the
difference visible and testable before the next one could rely on it.
